\begin{resumo}
O presente documento formaliza a Especificação de Requisitos de Software (SRS) do SIGEA-GTT (Sistema Inteligente de Gestão de Eventos Adversos -- \textit{Global Trigger Tool}), uma plataforma computacional hospitalar voltada à auditoria retrospectiva de prontuários, ensino clínico e rastreamento de eventos adversos, desenvolvida no âmbito do Departamento de Computação (DCOMP) e do Departamento de Enfermagem da Universidade Federal de Sergipe (UFS). Em conformidade com os preceitos normativos da ISO/IEC/IEEE 29148:2018 e do modelo de qualidade de produto ISO/IEC 25010, este documento especifica exaustivamente 25 Requisitos Funcionais (RF01 a RF25) e 10 Requisitos Não Funcionais (RNF01 a RNF10). Os requisitos funcionais delimitam os módulos de autenticação institucional restrita, parametrização do catálogo de módulos e gatilhos clínicos do \textit{Institute for Healthcare Improvement} (IHI), classificação de gravidade de dano NCC MERP, gestão de turmas acadêmicas, simulação de prontuários clínicos em formato JSONB, auditoria retrospectiva sob cronômetro estrito de 20 minutos, aplicação integrada de seis Ferramentas da Qualidade (Diagrama de Ishikawa 6M, Matriz GUT com cálculo dinâmico, Matriz 5W2H, Ciclo PDCA, Matriz SWOT e Painel de Ideação) e consolidação de indicadores epidemiológicos hospitalares e de proficiência pedagógica. Os requisitos não funcionais estabelecem restrições rigorosas de ergonomia visual desktop para resoluções homologadas (WUXGA, HD+, HD e Ultrawide 21:9), eficiência de resposta da API REST inferior a 200 milissegundos, segurança criptográfica com autenticação \textit{stateless} JWT e algoritmo BCrypt, isolamento total para conformidade com a Lei Geral de Proteção de Dados (LGPD) e critérios inegociáveis de \textit{Quality Gate} com 100\% de cobertura de testes automatizados e 100\% de documentação técnica de código. O documento inclui ainda o diagrama de casos de uso em UML 2.5 e a matriz de rastreabilidade bidirecional de requisitos.

\vspace{\onelineskip}
\noindent
\textbf{Palavras-chave}: Engenharia de Requisitos. Especificação de Software. Global Trigger Tool. Segurança do Paciente. Sistemas de Informação em Saúde. Universidade Federal de Sergipe.
\end{resumo}
